\documentclass[11pt]{article}
\usepackage[a4paper, top=18mm, bottom=18mm, left=20mm, right=20mm]{geometry}
\usepackage[T2A]{fontenc}
\usepackage[utf8]{inputenc}
\usepackage[ukrainian]{babel}
\usepackage{graphicx}
\usepackage[export]{adjustbox}
\usepackage{amsmath}
\usepackage{amssymb}
\graphicspath{{/app/Quiz/Scripts/2023/ProgAll/15BinTree/}}
\setlength{\parindent}{0pt}
\setlength{\parskip}{6pt}
\emergencystretch=3em
\pagestyle{empty}
\begin{document}
%begin
{\centering\Large\bfseries Контрольна робота\par}
\medskip
\noindent\rule{\linewidth}{0.5pt}
\smallskip
\noindent\textbf{Варіант \No\,4}\hfill \textit{ПІБ:}\,\underline{\hspace{60mm}}
\vspace{4mm}

%beginitem
1. \textbf{Що буде виведено за виконання фрагменту коду?}
d=['0','1','2','3','4']; d.insert(1,[1,0]); print(d)\par\vspace{5mm}
%enditem

%beginitem
2. 
Записати ВИРАЗ, значенням якого є список, що складається з кубів цілих чисел від 18 до 2008 (включно), що не діляться на жодне з чисел 2, 3, записаних в оберненому порядку.\par\vspace{5mm}
%enditem

%beginitem
3. 
Написати програму, що запитує у користувача значення дійсних параметрів $x$ та $y$, обчислює для них значення виразу 
$$\frac{\sqrt x - 59}{(\ctg x - 0.9) (y - 66)}$$ 
та повiдомляє введенi данi та результат обчислення.
 Оцінюється правильність та якість коду.

\par\vspace{5mm}
%enditem

%beginitem
4. 
Написати програму, що запитує у користувача значення дійсних параметрів $x$ та $y$, обчислює для них значення виразу 
$$\frac{\sqrt x - 59}{(\ctg x - 0.9) (y - 66)}$$ 
та повiдомляє введенi данi та результат обчислення.

\par\vspace{5mm}
%enditem

%beginitem
5. 
Записати ВИРАЗ, значенням якого є множина, що складається з цілих чисел від 18 до 2008 (включно), що діляться на 7.\par\vspace{5mm}
%enditem

%beginitem
6. \textbf{Що буде виведено за виконання фрагменту коду?}

%code
6//6%6*6
%code
\par\vspace{5mm}
%enditem

%beginitem
7. \textbf{Що буде виведено за виконання фрагменту коду?}
%Оригинал был утерян. Требует отладки!
%code
__d = 0

class D:
    __d = 1
    
    def __init__(self):
        self.__d = 2
        __d = 3
        D.__d = 4

obj = D()
D.__d = 5
try:
    print(__d, end=' ')
    print(obj.__d, end=' ')
    print(D.__d, end=' ')
    print(obj.__d, end=' ')
    print(obj._D__d)
except:
    print('ERR')
%code
\par\vspace{5mm}
%enditem

%beginitem
8. \textbf{Що буде виведено за виконання фрагменту коду?}
try:
    print(3, end=';')
    print(int('d4'), end=';')
    print(8, end=';')
except ZeroDivisionError:
    print(6, end=';')
except KeyboardInterrupt:
    print(7, end=';')
except:
    print('ERR', end=';')
else:
    print(1, end=';')
finally:
    print(2)
\par\vspace{5mm}
%enditem

%beginitem
9. \textbf{Що буде виведено за виконання фрагменту коду?}

print('R', end=' ')
for d in range(-7,-7,2):
    if d>=6:
        break
        print(d, end=' ')
    if d>=6:
        break
    else:
        print(d, end=' ')
    print(d)
\par\vspace{5mm}
%enditem

%beginitem
10. %goodpath:Scripts/2018/Mod2018_1/01-good.txt
\textbf{Відмітити все, що є однією правильною лексемою:}
( 1 ) +>

( 2 ) error

( 3 ) a2f

( 4 ) 0,5
\par\vspace{5mm}
%enditem

%end
\newpage
\end{document}

